% Generated draft: requires mathematical and source review.
\documentclass[11pt]{article}
\usepackage[margin=25mm]{geometry}
\usepackage{fontspec}
\setmainfont{Latin Modern Roman}
\usepackage{amsmath,amssymb,mathtools,unicode-math}
\setmathfont{Latin Modern Math}
\usepackage{enumitem,needspace,xurl,hyperref}
\setlength{\parindent}{0pt}
\setlength{\parskip}{5pt}
\setlength{\emergencystretch}{3em}
\newcommand{\N}{\mathbb{N}}
\newcommand{\Z}{\mathbb{Z}}
\newcommand{\Q}{\mathbb{Q}}
\newcommand{\R}{\mathbb{R}}
\newcommand{\C}{\mathbb{C}}
\newcommand{\F}{\mathbb{F}}
\newcommand{\A}{\mathbb{A}}
\newcommand{\PP}{\mathbb{P}}
\newcommand{\E}{\mathbb{E}}
\newcommand{\eps}{\varepsilon}
\newcommand{\dd}{\,\mathrm{d}}
\begin{document}
\section*{UnsolvedMath ID: 800375}
\textbf{Source problem code:} catalogue entry
\subsection{Definitions and mathematical statement}
Let \((W,S)\) be a Coxeter system, and let \(\ell:W\to\Z_{\geq 0}\) denote its length function. The (strong) Bruhat order is defined by \(x\leq y\) if a reduced expression for \(y\) contains a subword that is a reduced expression for \(x\). For \(u\leq v\), the Bruhat interval is the finite ranked poset
\[
[u,v]=\{x\in W:u\leq x\leq v\},
\]
with rank function \(x\mapsto\ell(x)-\ell(u)\). Put \(d=\ell(v)-\ell(u)\), and define its rank generating polynomial by
\[
R_{u,v}(q)=\sum_{x\in[u,v]}q^{\ell(x)-\ell(u)}
=\sum_{i=0}^{d}a_iq^i,
\qquad
a_i=\bigl|\{x\in[u,v]:\ell(x)-\ell(u)=i\}\bigr|.
\]
The interval is \emph{rank log-concave} when \(a_i^2\geq a_{i-1}a_{i+1}\) for every \(1\leq i\leq d-1\).

\textbf{Conjecture 13 (Brenti, 2024).} If \(W\) is a Weyl group, then every Bruhat interval in \(W\) is rank log-concave. Explicitly, for every \(u,v\in W\) with \(u\leq v\),
\[
\bigl|\{x\in[u,v]:\ell(x)-\ell(u)=i\}\bigr|^2
\geq
\bigl|\{x\in[u,v]:\ell(x)-\ell(u)=i-1\}\bigr|
\bigl|\{x\in[u,v]:\ell(x)-\ell(u)=i+1\}\bigr|
\quad(1\leq i\leq d-1).
\]
For \(d\leq1\), the condition is vacuous.

The source discusses rank generating functions for arbitrary Coxeter systems as background, but states this conjecture for Weyl groups. The usual finite crystallographic interpretation of \emph{Weyl group} is used here; the supplied excerpt does not specify any broader intended scope. The restriction \(u\leq v\) makes explicit the comparability implicit in the source's use of an interval. The supplied registry labels this conjecture open; that label is not an independently verified literature assessment.
\subsection{Short English statement}
Brenti's Conjecture 13 asserts that, for every comparable pair in a Weyl group, the numbers of elements at successive ranks of the Bruhat interval form a log-concave sequence. The draft uses the usual finite crystallographic meaning of Weyl group and requires human mathematical review.
\subsection{Sources}
\begin{itemize}
\item Local UnsolvedMath registry, permanent ID \texttt{800375}, exact source problem code \texttt{catalogue entry}, title \emph{Coxeter Groups and Unimodality --- Conjecture 13}. The supplied record attributes the conjecture to Brenti (2024) and labels it open.
\item Brenti, \emph{Some Open Problems on Coxeter Groups and Unimodality} (2024), Conjecture 13. Supplied source URL: \url{https://arxiv.org/abs/2410.09897}. The record reports access on 2026-10-16 and extraction by \texttt{arxiv\_latex}; these are supplied provenance details, not a claim of independent source access or verification.
\item The supplied background cites \texttt{[BjoEke]} for Bruhat intervals and their rank generating functions. No corresponding bibliographic entry is supplied, so this reference cannot be identified precisely from the record.
\end{itemize}
\end{document}
